\chapter{Superspace, Superfields, and Local Actions}
\label{ch:susy-superspace-superfields-local-actions}

\ConventionsLorentzian

Superspace is a supermanifold, understood as a locally super-ringed space.
Its odd coordinates are algebraic generators of local structure sheaves, not
points of an ordinary set.  This distinction is essential for a precise
relation between fermionic path-integral variables, component fields, and
operator-valued fermion fields.

\section{Superfields as Field-Variable Data}

\begin{definition}[Off-shell superfield multiplet]
\label{def:off-shell-superfield-multiplet}
An off-shell superfield multiplet is a finite collection of superfields and
constraints on them such that the supertranslation derivations act without
using Euler--Lagrange equations.  Its component fields are coordinates on a
space of field variables.  They are not, by definition, particle states in a
Hilbert space.
\end{definition}

If an action is supplied, the Euler--Lagrange equations may remove auxiliary
fields or impose differential constraints.  If gauge symmetry is present, the
physical field configurations are further quotients or cohomology classes.
Only after these operations and a Hilbert-space or Euclidean reconstruction
has been specified can one compare the resulting one-particle spectrum with
the particle representations of
Chapter~\ref{ch:susy-algebras-representations}.  This chapter therefore uses
the word \emph{multiplet} only for field-variable packages unless explicitly
stated otherwise.

\section{Affine Superspace}

Four-dimensional \(\mathcal N=1\) affine superspace has underlying ordinary
manifold \(\mathbb M^4\) and structure sheaf
\[
  \mathcal O_{\mathbb M^{4|4}}(U)
  =
  C^\infty(U)\otimes
  \Lambda(\theta^1,\theta^2,\bar\theta^{\dot1},\bar\theta^{\dot2})
\]
for open \(U\subset\mathbb M^4\).  A superfield is a section of this sheaf
or of a vector bundle tensored with it.  In local coordinates,
\[
  F(x,\theta,\bar\theta)
  =
  f(x)+\theta^\alpha f_\alpha(x)
  +\bar\theta_{\dot\alpha}\bar f^{\dot\alpha}(x)
  +\cdots+\theta^2\bar\theta^2 f_{\rm top}(x),
\]
where the expansion terminates because the odd coordinate algebra is finite.

The supertranslation generators act as derivations of the structure sheaf:
\[
  Q_\alpha
  =
  \frac{\partial}{\partial\theta^\alpha}
  -\ii\sigma^\mu_{\alpha\dot\alpha}\bar\theta^{\dot\alpha}\partial_\mu,
  \qquad
  \bar Q_{\dot\alpha}
  =
  -\frac{\partial}{\partial\bar\theta^{\dot\alpha}}
  +\ii\theta^\alpha\sigma^\mu_{\alpha\dot\alpha}\partial_\mu .
\]
The covariant derivatives are
\[
  D_\alpha
  =
  \frac{\partial}{\partial\theta^\alpha}
  +\ii\sigma^\mu_{\alpha\dot\alpha}\bar\theta^{\dot\alpha}\partial_\mu,
  \qquad
  \bar D_{\dot\alpha}
  =
  -\frac{\partial}{\partial\bar\theta^{\dot\alpha}}
  -\ii\theta^\alpha\sigma^\mu_{\alpha\dot\alpha}\partial_\mu .
\]
They anticommute with the supercharges and obey
\[
  \{D_\alpha,\bar D_{\dot\beta}\}
  =
  -2\ii\sigma^\mu_{\alpha\dot\beta}\partial_\mu .
\]

\section{Chiral Superfields}

\begin{definition}[Chiral superfield]
\label{def:chiral-superfield}
A chiral superfield is a superfield \(\Phi\) satisfying
\[
  \bar D_{\dot\alpha}\Phi=0 .
\]
Equivalently, in chiral coordinates
\[
  y^\mu=x^\mu+\ii\theta\sigma^\mu\bar\theta,
\]
it has the component expansion
\[
  \Phi(y,\theta)
  =
  \phi(y)+\sqrt2\,\theta^\alpha\psi_\alpha(y)+\theta^2F(y).
\]
\end{definition}

Here \(\phi\) is an even scalar field, \(\psi_\alpha\) an odd Weyl spinor
field, and \(F\) an even auxiliary field.  The component supersymmetry
transformations follow by acting with
\(\epsilon^\alpha Q_\alpha+\bar\epsilon_{\dot\alpha}\bar Q^{\dot\alpha}\)
and re-expanding in \(\theta\):
\[
  \delta\phi=\sqrt2\,\epsilon\psi,
\]
\[
  \delta\psi_\alpha
  =
  \sqrt2\,\epsilon_\alpha F
  +\ii\sqrt2\,\sigma^\mu_{\alpha\dot\alpha}\bar\epsilon^{\dot\alpha}
  \partial_\mu\phi,
\]
\[
  \delta F
  =
  \ii\sqrt2\,\bar\epsilon_{\dot\alpha}
  \bar\sigma^{\mu\dot\alpha\alpha}\partial_\mu\psi_\alpha .
\]
The auxiliary field makes the algebra close on components before using
Euler--Lagrange equations.

\section{Berezin Integration and Local Superspace Actions}

Berezin integration is the projection functional onto the top odd coefficient
of the exterior algebra:
\[
  \int d^2\theta\,\theta^2=1,\qquad
  \int d^2\bar\theta\,\bar\theta^2=1.
\]
As a linear functional on the odd coordinate algebra, it is the integration
operation intrinsic to the locally super-ringed-space description.  A local
superspace action for chiral fields has the form
\[
  S
  =
  \int d^4x\,d^4\theta\,K(\Phi^i,\bar\Phi^{\bar j})
  +
  \left[
  \int d^4x\,d^2\theta\,W(\Phi^i)+{\rm h.c.}
  \right],
\]
where \(K\) is a real function of superfields and \(W\) is holomorphic in the
chiral superfields.  The \(d^4\theta\) term is a \(D\)-term; the
\(d^2\theta\) term is an \(F\)-term.

For one chiral field with \(K=\bar\Phi\Phi\), the component expansion gives
\[
  \mathcal L
  =
  -\partial_\mu\bar\phi\,\partial^\mu\phi
  -\ii\bar\psi\bar\sigma^\mu\partial_\mu\psi
  +\bar F F
  +
  \left[
  F W'(\phi)
  -\frac{1}{2}W''(\phi)\psi^\alpha\psi_\alpha
  +{\rm h.c.}
  \right]
\]
in the mostly-plus convention.  Eliminating \(F\) by its algebraic
Euler--Lagrange equation gives \(F=-\overline{W'(\phi)}\) and potential
\[
  V(\phi,\bar\phi)=|W'(\phi)|^2 .
\]

\section{Operator Fields and Grassmann Variables}

Fermionic component fields in a Hilbert-space QFT are operator-valued
distributions satisfying graded locality.  Fermionic variables in a
superspace or path-integral expression are elements of odd pieces of a
Grassmann algebra.  The bridge between them is the correlation functional:
Grassmann sources and Berezin integration encode antisymmetric multilinear
forms that represent vacuum or Euclidean correlation functions of the
operator-valued fermion fields after reconstruction or perturbative
definition.

\section{Supersymmetric Wilsonian Schemes}

The later dynamics of four-dimensional \(\mathcal N=1\) and \(\mathcal N=2\)
gauge theories will use Wilsonian effective actions.  Such statements require
a scheme, not only a formal superspace expression.

\begin{definition}[Supersymmetric Wilsonian scheme]
\label{def:supersymmetric-wilsonian-scheme}
A supersymmetric Wilsonian scheme on a scale interval
\(\Lambda_{\rm IR}<\Lambda<\Lambda_{\rm UV}\) consists of the following data:
\begin{enumerate}[label=(\roman*)]
\item a regulated space of field variables \(\mathfrak F_\Lambda\), including
ordinary fields, superfields, ghosts, antifields, and background sources when
they are part of the construction;
\item an integration prescription with Lagrangian functional \(S_\Lambda\):
for gauge theories this means either a BV finite-cutoff integral with
bracket, Laplacian, half-density, and master action specified, or a lattice
construction with gauge-invariant link-variable data and blocking maps;
\item odd derivations or Ward operators representing the supersymmetry
algebra on the regulated variables, with a stated closure relation: off shell,
on shell, modulo gauge transformations, or modulo the BV differential;
\item a coarse-graining map from scale \(\Lambda_2\) to scale \(\Lambda_1\)
for \(\Lambda_{\rm IR}<\Lambda_1<\Lambda_2<\Lambda_{\rm UV}\), satisfying a
composition law up to finite local coordinate changes;
\item a coordinate chart on local and quasilocal functionals, including the
separation of \(D\)-terms, \(F\)-terms, gauge kinetic terms, theta terms,
and local counterterms allowed by the regulator;
\item a prescription for composite operator insertions and background-field
variations, stated separately from the regularization of the functional
integral itself.
\end{enumerate}
The scheme is called supersymmetric if the regulated Ward identity is exact
at finite \(\Lambda\), or if the failure is a local breaking functional with a
specified restoration prescription and anomaly class.
\end{definition}

Holomorphic or nonrenormalization statements are statements about coordinates
in such a scheme.  They are not statements about an undefined path-integral
symbol, nor are they automatically statements about the one-particle
irreducible effective action.  A physical statement extracted from a
Wilsonian action must be invariant under the finite coordinate changes allowed
by the scheme, or the coordinate dependence must be part of the statement.
This distinction is essential in the exact dynamics of four-dimensional
supersymmetric gauge theories.

There is a second distinction which will be used repeatedly.  The BV
formalism can encode gauge symmetry and supersymmetry on a space of fields,
ghosts, antifields, and background superfields, and it can express a
symmetry-preserving coarse-graining step as a canonical or quantum-canonical
operation once a regulator has been supplied.  This does not by itself produce
an ultraviolet regularization.  A regulated BV Laplacian, a finite
field-variable space, a cutoff propagator, an integration cycle, and a proof
that the regulated quantum master equation and Ward identities hold are
additional data.  Thus an off-shell-superfield BV description is a natural
language for supersymmetric Wilsonian integration, but it is not identical to
a manifestly supersymmetric regulator.

For supersymmetric gauge theories, the Wilsonian effective theory is therefore
formulated through BV data unless a lattice construction is explicitly
available.  The BV datum must include the regulated odd symplectic structure,
normal half-density, BV Laplacian, integration cycle, master action, and
coarse-graining map.  A lattice datum must include the finite gauge-field
configuration space, fermion and auxiliary-field regulator, exact gauge
symmetry, reflection or Hilbert-space positivity status when needed, and the
blocking map for observables.  Holomorphic coordinates, moduli-space
coordinates, and exact superpotential arguments are certified only relative to
one of these gauge-compatible finite-cutoff frameworks.

For four-dimensional supersymmetric gauge theories we shall not assume the
existence of a completely satisfactory manifestly supersymmetric
regularization scheme.  Dimensional reduction is a perturbative prescription:
loop momenta are analytically continued away from four dimensions while
spinor and gauge algebra are partly kept four-dimensional.  A calculation
performed with this prescription is therefore certified only after the
algebraic rules used at the relevant loop order have been stated and checked,
including identities involving \(\epsilon_{\mu\nu\rho\sigma}\), chiral
matrices, and evanescent tensor structures.  Higher-covariant-derivative or
heat-kernel regulators supplemented by Pauli--Villars-type fields are likewise
candidate schemes, not axioms: one must state the regulator fields, statistics,
masses, gauge and supersymmetry transformations, BV pairing, decoupling limit,
local counterterm freedom, and anomaly class before using such a construction
as a proved input.

\begin{openproblem}[Regulators preserving supersymmetry]
\label{op:susy-regulators-preserving-supersymmetry}
Classify regulator schemes for supersymmetric QFT by the exact symmetries
they preserve, the Ward identities they realize before removing the
regulator, and the local counterterms required to restore the remaining
identities.  Dimensional reduction must be treated as a formal perturbative
scheme whose consistency is checked at the loop order being used, rather than
as a mathematical definition of a supersymmetric continuum theory.
Higher-derivative heat-kernel schemes, Pauli--Villars completions, lattice
constructions, and BV-compatible exact-RG schemes must be compared at the
level of their regulated master equations or finite lattice Ward identities,
locality of counterterms, decoupling limits, and anomaly cohomology before
nonrenormalization or anomaly statements are used.
\end{openproblem}

\begin{openproblem}[Wilsonian foundations for four-dimensional supersymmetric gauge dynamics]
\label{op:susy-wilsonian-foundations-gauge-dynamics}
Develop the exact dynamics of four-dimensional \(\mathcal N=1\) and
\(\mathcal N=2\) gauge theories inside explicitly defined supersymmetric
Wilsonian schemes.  The target is to formulate holomorphic couplings,
effective superpotentials, quantum moduli spaces, Seiberg-Witten effective
actions, anomaly matching, decoupling, and duality tests as statements about
specified Wilsonian coordinates and scheme-invariant observables, with the
relation to 1PI effective actions stated separately.
\end{openproblem}

\begin{openproblem}[Supersymmetric examples across dimensions]
\label{op:susy-examples-across-dimensions}
Develop the main classes of supersymmetric QFTs in dimensions other than four
with the same separation of algebra, field-variable realization, regulator,
Wilsonian coordinate, and observable.  In two dimensions this includes
\(\mathcal N=(2,2)\) Landau--Ginzburg models, nonlinear sigma models with
Calabi--Yau target, and gauged linear sigma models, where the relation among
UV superpotential data, Kähler data, chiral rings, RG flow, and infrared CFT
must be stated as a theorem, construction, controlled argument, or open
problem rather than as a phase slogan.  In three dimensions this includes
supersymmetric Chern--Simons--matter theories, where level quantization,
framing, parity anomalies, boundary conditions, monopole operators, and the
meaning of Wilsonian coordinates for topological terms must be part of the
definition.  In six dimensions this includes superconformal theories for
which a conventional Lagrangian may not be part of the data; the tensor-branch
effective theory, anomaly polynomial, BPS strings, compactifications, and
extended operators then become diagnostic objects rather than replacements
for a nonperturbative definition.
\end{openproblem}

\begin{definition}[Localization datum]
\label{def:susy-localization-datum}
A localization datum for a supersymmetric functional integral consists of:
\begin{enumerate}[label=(\roman*)]
\item a regulated integration space or integration cycle, including its
orientation, boundary, singular strata, and gauge quotient if present;
\item an odd symmetry \(Q\) acting on the regulated variables, together with
the precise bosonic transformation represented by \(Q^2\);
\item a \(Q\)-invariant action and \(Q\)-closed observables;
\item a deformation \(t\,Q V\) whose convergence, boundary behavior, and
large-\(t\) limit are controlled;
\item the fixed locus of \(Q\), its normal complex, zero modes, and determinant
or Euler-class prescription;
\item a contour or residue prescription for remaining finite-dimensional
integrals, including the rule used at singular hyperplane arrangements.
\end{enumerate}
\end{definition}

Finite-dimensional equivariant localization is a theorem only after the
geometric hypotheses of the theorem have been stated.  An infinite-dimensional
supersymmetric path integral inherits no automatic theorem from this analogy.
One must specify the regulator or finite-dimensional approximation, prove or
assume the absence of boundary contributions, and describe singular loci.  In
particular, Jeffrey--Kirwan residues are contour prescriptions tied to
hyperplane arrangements and chamber choices; they are not a substitute for a
definition of the original path integral.  Likewise, zero-size gauge
instantons are points on a singular compactification or boundary of a moduli
space; whether they are legitimate integration objects depends on the chosen
compactification, measure, and local QFT interpretation.

In finite-dimensional reductions of gauge theories, a Jeffrey--Kirwan residue
typically appears after abelianization or gauge fixing has replaced part of
the integral by a meromorphic form on the complexified Cartan algebra with
hyperplane divisors from charged fluctuations.  The chamber vector encodes
orientation data such as a stability parameter, Fayet--Iliopoulos parameter,
or contour at infinity.  A QFT localization argument must therefore derive the
meromorphic form and this chamber data from the regulated integration cycle;
otherwise the residue is an additional prescription selecting one boundary
value of a singular integral.

\begin{openproblem}[Rigorous supersymmetric localization]
\label{op:rigorous-supersymmetric-localization}
Develop supersymmetric localization in QFT from regulated localization data.
The target is to identify when the localized integral is well-defined, when
the deformation argument has no boundary contribution, how singular instanton
strata and noncompact Coulomb or Higgs directions are compactified or
excluded, and how residue prescriptions such as Jeffrey--Kirwan arise from
the declared integration cycle rather than from a formal rule.
\end{openproblem}
